\chapter{Methods of Solving First-Order ODEs}

% -------------------------------------------------------
\section{Separable Equations}
% -------------------------------------------------------

\begin{definition}[Separable Equation]
A first-order ODE is \textbf{separable} if it can be written in the form
\[
  \frac{dy}{dx} = f(x)\,g(y).
\]
\end{definition}

\subsection{Method}

Separate the variables so that all $y$-terms are on one side and all $x$-terms on the other:
\[
  \frac{dy}{g(y)} = f(x)\,dx.
\]
Integrate both sides:
\[
  \int \frac{dy}{g(y)} = \int f(x)\,dx + C.
\]

\begin{example}
Solve $\displaystyle x\,\frac{dy}{dx} + y = 0$.

\begin{solution}
Rewrite:
\[
  \frac{dy}{dx} = -\frac{y}{x}.
\]
Separate:
\[
  \frac{dy}{y} = -\frac{dx}{x}.
\]
Integrate both sides:
\[
  \ln|y| = -\ln|x| + C_1.
\]
Exponentiate:
\[
  |y| = e^{C_1}\cdot\frac{1}{|x|}.
\]
Writing $C = \pm e^{C_1}$ as an arbitrary constant:
\[
  \boxed{y = \frac{C}{x}.}
\]
\end{solution}
\end{example}

% -------------------------------------------------------
\section{Linear First-Order Equations}
% -------------------------------------------------------

\begin{definition}[Standard Linear Form]
A first-order ODE is in \textbf{standard linear form} when written as
\[
  \frac{dy}{dx} + P(x)\,y = Q(x).
\]
\end{definition}

\subsection{Method: Integrating Factor}

\begin{theorem}[Integrating Factor Method]
Multiply both sides of $y' + P(x)y = Q(x)$ by the \textbf{integrating factor}
\[
  \IF(x) = e^{\int P(x)\,dx}.
\]
The left-hand side becomes an exact derivative:
\[
  \frac{d}{dx}\!\left[\IF(x)\cdot y\right] = \IF(x)\cdot Q(x).
\]
Integrate both sides to obtain $y$.
\end{theorem}

\begin{example}
Solve $\displaystyle\frac{dy}{dx} + 3y = e^{2x}$.

\begin{solution}
Here $P(x)=3$ and $Q(x)=e^{2x}$.

\textbf{Step 1 — Integrating factor:}
\[
  \IF = e^{\int 3\,dx} = e^{3x}.
\]

\textbf{Step 2 — Multiply through:}
\[
  e^{3x}\frac{dy}{dx} + 3e^{3x}y = e^{3x}\cdot e^{2x} = e^{5x}.
\]
The left-hand side is $\dfrac{d}{dx}\!\left(ye^{3x}\right)$, so:
\[
  \frac{d}{dx}\!\left(ye^{3x}\right) = e^{5x}.
\]

\textbf{Step 3 — Integrate:}
\[
  ye^{3x} = \int e^{5x}\,dx = \frac{e^{5x}}{5} + C.
\]

\textbf{Step 4 — Solve for $y$:}
\[
  \boxed{y = \frac{e^{2x}}{5} + Ce^{-3x}.}
\]
\end{solution}
\end{example}

% -------------------------------------------------------
\section{Homogeneous First-Order Equations}
% -------------------------------------------------------

\begin{definition}[Homogeneous First-Order ODE]
A first-order ODE is \textbf{homogeneous} (in the function-form sense) if it can be written as
\[
  \frac{dy}{dx} = f\!\left(\frac{y}{x}\right).
\]
\end{definition}

\subsection{Method: Substitution}

Let $v = \dfrac{y}{x}$, so $y = vx$. Then
\[
  \frac{dy}{dx} = v + x\,\frac{dv}{dx}.
\]
Substituting converts the equation into a separable ODE in $v$ and $x$.

\begin{remark}
After solving for $v(x)$, substitute back using $y = vx$.
\end{remark}

% -------------------------------------------------------
\section{Exact Equations}
% -------------------------------------------------------

\begin{definition}[Exact Equation]
The equation
\[
  M(x,y)\,dx + N(x,y)\,dy = 0
\]
is \textbf{exact} if there exists a function $F(x,y)$ such that
\[
  \frac{\partial F}{\partial x} = M \qquad\text{and}\qquad \frac{\partial F}{\partial y} = N,
\]
in which case the general solution is $F(x,y) = C$.
\end{definition}

\begin{theorem}[Exactness Condition]
The equation $M\,dx + N\,dy = 0$ is exact if and only if
\[
  \frac{\partial M}{\partial y} = \frac{\partial N}{\partial x}.
\]
\end{theorem}

\subsection{Method}

\begin{enumerate}
  \item Verify $\partial M/\partial y = \partial N/\partial x$.
  \item Integrate $M$ with respect to $x$: $F = \int M\,dx + g(y)$.
  \item Differentiate $F$ with respect to $y$ and set equal to $N$ to find $g'(y)$.
  \item Integrate to obtain $g(y)$; the solution is $F(x,y) = C$.
\end{enumerate}
